\noindent\textbf{Verbatim problem statement (from file).}
\emph{``Let $h(n)$ count the number of incongruent sets of $n$ points in $\R^2$ which minimise the diameter subject to the constraint that $d(x,y)\ge 1$ for all points $x\ne y$. Is it true that $h(n)\to\infty$? It is not even known whether $h(n)\ge 2$ for all large $n$.''}

\medskip
\noindent\textbf{Definitions and conventions.}
For $n\ge 1$, consider the class $\mathcal{C}_n$ of all $n$-point sets $A\subset\R^2$ with pairwise distances at least $1$.
Define the minimal achievable diameter
\[\Delta_n := \inf\{\diam(A): A\in\mathcal{C}_n\}.
\]
Call $A\in\mathcal{C}_n$ \emph{diameter-minimising} if $\diam(A)=\Delta_n$.
Two sets are \emph{congruent} if one can be mapped to the other by an isometry of $\R^2$.
Define $h(n)$ as the number of congruence classes of diameter-minimising sets in $\mathcal{C}_n$.

\medskip
\noindent\textbf{Question.}
Does $h(n)\to\infty$ as $n\to\infty$?
